\chapter{Monte Carlo}
\label{ch:mc}

\section{The sample average}

\begin{defbox}
If $x^{(1)},\ldots,x^{(N)}$ are drawn independently from $p$, the Monte Carlo estimator of $\E_p[f]$ is
\[
\hat f_N=\frac1N\sum_{n=1}^N f(x^{(n)}).
\]
It is unbiased. If $\mathrm{Var}_p(f)=\sigma_f^2<\infty$, then $\mathrm{Var}(\hat f_N)=\sigma_f^2/N$.
\end{defbox}

\begin{intbox}
Unbiased means the average of the estimator, over datasets of samples, equals the integral. It does not mean one sample is exact. The error typically shrinks like $1/\sqrt{N}$, because the standard deviation is $\sigma_f/\sqrt{N}$. Lecture~2, slides~34--35, and Homework~1, problem~2, turn that variance into a sample size through Chebyshev or a Chernoff bound. Chebyshev uses only the variance. A Chernoff bound uses the moment generating function and can demand fewer samples while remaining a guarantee.
\end{intbox}

\begin{resultbox}[Chebyshev sample size]
If $\sigma_f^2<\infty$, then for every $\varepsilon>0$,
\[
\mathbb P\bigl(|\hat f_N-\E_p f|\ge\varepsilon\bigr)\le\frac{\sigma_f^2}{N\varepsilon^2}.
\]
The right-hand side tends to $0$ (weak law). The sample size needed for deviation $\varepsilon$ and failure probability $\delta$ is $N\ge\sigma_f^2/(\delta\varepsilon^2)$. Root-mean-square error is $O(N^{-1/2})$.
\end{resultbox}

A Chernoff bound replaces the variance by the moment generating function. For a $\chi^2_N$ law the mgf is $(1-2\lambda)^{-N/2}$ for $\lambda<1/2$, and optimizing Markov's inequality on each tail produces explicit rate functions $I_\pm(t)$. The resulting $N$ is still a finite-sample guarantee, typically tighter than Chebyshev and looser than a CLT quantile (which is only asymptotic).

\section{Importance sampling}

When draws from $p$ are unavailable, draw $x\sim q$ instead and correct the weight:
\[
\E_p[f]=\E_q\Bigl[f(x)\frac{p(x)}{q(x)}\Bigr],
\]
provided $q$ is positive wherever $p f$ is nonzero. The weight $p/q$ is the importance weight. The estimator is unbiased for any such $q$. Its variance is small when $q$ is large where $\lvert f\rvert p$ is large, and it is huge when $q$ is tiny on that set.

\begin{resultbox}[Second-moment criterion]
The one-sample variance $\mathrm{Var}_q(fp/q)$ is finite if and only if $\E_q[(fp/q)^2]<\infty$. Among nonnegative integrands, the zero-variance proposal is $q^\star(x)\propto f(x)p(x)$, which requires the unknown integral as a normalizer and is therefore not a practical proposal.
\end{resultbox}

For the scalar example $f(x)=x^2$, $p=\N(0,1)$, and Gaussian proposals $q=\N(0,s^2)$, the estimator is unbiased for every $s>0$, the second moment is finite if and only if $s>1/\sqrt 2$, and the variance-minimizing scale among Gaussians is $s=\sqrt 3$.

\begin{pitbox}
Lecture~4, slides~23--29, shows a proposal that misses the posterior mass. The likelihood estimator then has rare enormous weights. The ELBO replaces that high-variance estimator by a bound whose gap is exactly a KL. You give up unbiasedness of the likelihood and gain a gradient you can train.
\end{pitbox}

\section{Perplexity}

For a discrete model, perplexity is $\exp\bigl(-\E_{x\sim p_{\mathrm{data}}}\log p_\theta(x)\bigr)$, or the same expression with the average over a finite test set. It is an exponential of a cross-entropy. Write $\ell(x)=-\log p_\theta(x)$, $\mathrm{PPL}=\exp(\E[\ell])$, and $\widehat{\mathrm{PPL}}_N=\exp(\bar\ell_N)$ with $\bar\ell_N=N^{-1}\sum_n\ell(x_n)$.

\begin{resultbox}[Bias and consistency of plug-in perplexity]
Jensen for the convex map $\exp$ gives $\E[\widehat{\mathrm{PPL}}_N]\ge\mathrm{PPL}$ for every finite $N$, with equality if and only if $\ell$ is $p_{\mathrm{data}}$-a.s.\ constant. If $\E|\ell|<\infty$, the weak law and continuity of $\exp$ give consistency in probability. Under a third-moment assumption the bias admits the expansion
\[
\E[\widehat{\mathrm{PPL}}_N]=\mathrm{PPL}\Bigl(1+\frac{\mathrm{Var}(\ell)}{2N}+o(N^{-1})\Bigr).
\]
\end{resultbox}
When data are Gaussian and the model is $\N(0,\sigma^2)$ with variance ratio $r=\mathrm{Var}_{\mathrm{data}}/\sigma^2$, the plug-in has a closed form involving the chi-squared mgf, finite if and only if $N>r$, and the asymptotic coefficient matches $\mathrm{Var}(\ell)/2=r^2/4$.

Homework~1 asks for that expansion. The point to remember is that the log of an average is not the average of the logs---here the exponential of an average is biased upward.

\section{Homework 1 problems}

The KL and Jensen--Shannon parts of Homework~1 problem~2 live in Chapter~\ref{ch:info}. The Monte Carlo, importance sampling, and perplexity parts are collected here.

\begin{hwqbox}[Homework 1, problem 2.1]
Show that the Monte Carlo average $f_N=N^{-1}\sum_{n\le N}f(x_n)$ is unbiased when $\E_p|f|<\infty$, compute its variance when $\mathrm{Var}_p(f)=\sigma_f^2<\infty$, and deduce Chebyshev's bound and the weak law.
\end{hwqbox}
\begin{hwsolbox}[Homework 1, problem 2.1]
If $\E_p|f|<\infty$, then $\E[f_N]=\E_p[f]$, so $f_N$ is unbiased. If $\mathrm{Var}_p(f)=\sigma_f^2<\infty$, independence gives $\mathrm{Var}(f_N)=\sigma_f^2/N$.

Chebyshev: for every $\varepsilon>0$,
\[
\mathbb P\big(|f_N-\E_p f|\ge\varepsilon\big)\le\frac{\sigma_f^2}{N\varepsilon^2}.
\]
The right-hand side tends to $0$, which is the weak law, with rate $O(1/N)$ in probability deviation and $O(N^{-1/2})$ in root-mean-square error. The sample size needed for deviation $\varepsilon$ and probability $\delta$ is $N\ge\sigma_f^2/(\delta\varepsilon^2)$.
\end{hwsolbox}

\begin{hwqbox}[Homework 1, problem 2.1(a)]
Let $X\sim\N(1,4)$ and $f(X)=(X-1)^2$. Bound or approximate the $N$ needed so that $\mathbb P(|f_N-4|\ge 1/2)\le 0.05$ by (i) Chebyshev, (ii) CLT, and (iii) Chernoff on the underlying chi-squared. Round Chernoff to the nearest $50$, and compare the three answers.
\end{hwqbox}
\begin{hwsolbox}[Homework 1, problem 2.1(a)]
Let $Z=(X-1)/2\sim\N(0,1)$. Then $f(X)=4Z^2$, so $\E[f]=4$ and $\mathrm{Var}(f)=16\mathrm{Var}(Z^2)=16(3-1)=32$.

Each $(X_n-1)^2/4\sim\chi^2_1$, hence $f_N=(4/N)\chi^2_N$. The event $|f_N-4|\ge 1/2$ is $|\chi^2_N-N|\ge N/8$, i.e.\ a relative deviation $t=1/8$.

(i) Chebyshev. $\mathrm{Var}(f_N)=32/N$, so the probability is at most $32/(N\cdot 1/4)=128/N$. The bound is at most $0.05$ once $N\ge 128/0.05=2560$. This $N$ is guaranteed.

(ii) CLT. Use $\sqrt{32/N}$ as the standard deviation and the normal quantile $1.96$:
\[
\frac{1/2}{\sqrt{32/N}}\ge 1.96
\iff N\ge 32\cdot(3.92)^2=491.72\ldots
\]
So $N\ge 492$. This is an approximation, not a finite-sample guarantee.

(iii) Chernoff. The mgf of $\chi^2_N$ is $(1-2\lambda)^{-N/2}$ for $\lambda<1/2$. Optimizing the Markov bound on each tail produces, for $t\in(0,1)$,
\begin{align*}
\mathbb P(\chi^2_N\ge N(1+t))&\le\big((1+t)e^{-t}\big)^{N/2}=e^{-NI_+(t)},\\
\mathbb P(\chi^2_N\le N(1-t))&\le\big((1-t)e^{t}\big)^{N/2}=e^{-NI_-(t)},
\end{align*}
with $I_+(t)=\tfrac12(t-\log(1+t))$ and $I_-(t)=\tfrac12(-\log(1-t)-t)$. At $t=1/8$, $I_+\approx 0.0036085$ and $I_-\approx 0.0042657$. The sum of the two tails drops to $0.05$ or below at $N=950$. Nearest $50$ is $950$. This $N$ is guaranteed.

Chebyshev and Chernoff are theorems about this finite $N$; the CLT is only an asymptotic shape. Chebyshev is loose ($2560$). Chernoff is tighter ($950$) but still a bound. The CLT suggests about $492$ and is the most optimistic precisely because it is not guaranteed.
\end{hwsolbox}

\begin{hwqbox}[Homework 1, problem 2.1(b)]
Importance-sample $\E_p[x^2]=1$ for $p=\N(0,1)$ with proposals $q=\N(0,s^2)$. Show unbiasedness for every $s>0$, find when the variance is finite, minimize that variance over $s$, and identify the zero-variance proposal over all densities.
\end{hwqbox}
\begin{hwsolbox}[Homework 1, problem 2.1(b)]
\[
\E_q\Big[f(x)\frac{p(x)}{q(x)}\Big]=\int x^2 p(x)\,dx=1
\]
whenever $s>0$, because $q>0$ everywhere and $x^2p$ is integrable. So $\hat I_N$ is unbiased for every $s>0$, and $\mathrm{Var}(\hat I_N)=N^{-1}\mathrm{Var}_q(fp/q)$.

The second moment is finite if and only if $s>1/\sqrt 2$. Indeed
\[
\frac{p(x)^2}{q(x)}=s(2\pi)^{-1/2}\exp\big(-x^2(1-\tfrac1{2s^2})\big),
\]
and $\int x^4 p^2/q$ converges precisely when the quadratic rate $1-1/(2s^2)$ is positive. Evaluating the Gaussian moment gives
\[
\E_q\Big[\Big(x^2\frac pq\Big)^2\Big]=\frac{3s^6}{(2s^2-1)^{5/2}},\qquad
\mathrm{Var}_q\Big(x^2\frac pq\Big)=\frac{3s^6}{(2s^2-1)^{5/2}}-1.
\]
Minimize $v(u)=3u^3/(2u-1)^{5/2}$ for $u=s^2>1/2$. The logarithmic derivative $3/u-5/(2u-1)$ vanishes at $u=3$, and $v\to\infty$ at both endpoints, so the best Gaussian scale is $s=\sqrt 3$, with minimal one-sample variance $(81\sqrt 5-125)/125\approx 0.44897$. Ordinary Monte Carlo ($s=1$) has variance $2$.

Over all proposals, $f\ge 0$, so the zero-variance proposal is $q^\star(x)\propto f(x)p(x)=x^2 p(x)$. It is already normalized because $\E_p[x^2]=1$:
\[
q^\star(x)=x^2\cdot(2\pi)^{-1/2}e^{-x^2/2}.
\]
Then $fp/q^\star=1$ for $x\ne 0$, so the importance-sampling estimator has variance $0$. Using it requires the unknown integral as a normalizing constant.
\end{hwsolbox}

\begin{hwqbox}[Homework 1, problem 2.2]
Define $\widehat{\mathrm{PPL}}_N=\exp(\bar\ell_N)$ with $\bar\ell_N=N^{-1}\sum_{n\le N}\ell(x_n)$ and $\mathrm{PPL}=\exp(\E[\ell])$. Show that the estimator is biased upward for every finite $N$, characterize equality, and prove consistency when $\E|\ell|<\infty$.
\end{hwqbox}
\begin{hwsolbox}[Homework 1, problem 2.2]
Jensen's inequality for the convex map $\exp$ gives
\[
\E[\widehat{\mathrm{PPL}}_N]\ge\exp(\E[\bar\ell_N])=\exp(\E[\ell])=\mathrm{PPL},
\]
for every finite $N$. Equality holds if and only if $\bar\ell_N$ is almost surely constant, i.e.\ $\ell(x)$ is $p_{\mathrm{data}}$-a.s.\ constant. In particular $\widehat{\mathrm{PPL}}_N$ is not unbiased unless $\ell$ is a.s.\ constant.

If $\E|\ell|<\infty$, the weak law gives $\bar\ell_N\to\E[\ell]$ in probability, and continuity of $\exp$ gives $\widehat{\mathrm{PPL}}_N\to\mathrm{PPL}$ in probability. The estimator is consistent. The strong law upgrades this to almost-sure convergence.
\end{hwsolbox}

\begin{hwqbox}[Homework 1, problem 2.2(a)]
Assuming $\E|\ell|^3<\infty$, expand $\E[\widehat{\mathrm{PPL}}_N]/\mathrm{PPL}$ through order $1/N$ and identify the coefficient $c=\mathrm{Var}(\ell)/2$.
\end{hwqbox}
\begin{hwsolbox}[Homework 1, problem 2.2(a)]
Write $\delta_N=\bar\ell_N-\E[\ell]$, so $\E\delta_N=0$, $\E\delta_N^2=\mathrm{Var}(\ell)/N$, and $\E\delta_N^3=O(1/N^2)$. Then
\[
\E[e^{\delta_N}]=1+\tfrac12\E[\delta_N^2]+O(1/N^2)=1+\frac{\mathrm{Var}(\ell)}{2N}+o(N^{-1}).
\]
Hence $c=\mathrm{Var}(\ell)/2$.
\end{hwsolbox}

\begin{hwqbox}[Homework 1, problem 2.2(b)]
For Gaussian data of variance $s^2$ and model $\N(0,\sigma^2)$ with $r=s^2/\sigma^2$, derive the closed form of $\E[\widehat{\mathrm{PPL}}_N]$, the condition $N>r$ for finiteness, the asymptotic coefficient $c=r^2/4$, and the interpretation when $r>1$.
\end{hwqbox}
\begin{hwsolbox}[Homework 1, problem 2.2(b)]
\[
\ell(x)=\tfrac12\log(2\pi\sigma^2)+\frac{x^2}{2\sigma^2},\qquad
\E[\ell]=\tfrac12\log(2\pi\sigma^2)+\tfrac r2,
\]
so $\mathrm{PPL}=\sqrt{2\pi\sigma^2}\,e^{r/2}$. With $\sum_n x_n^2=s^2\chi^2_N$,
\[
\widehat{\mathrm{PPL}}_N=\sqrt{2\pi\sigma^2}\,\exp\big(\tfrac{r}{2N}\chi^2_N\big).
\]
The chi-squared mgf at $\lambda=r/(2N)$ exists if and only if $r/N<1$, i.e.\ $N>r$, and then
\[
\E[\widehat{\mathrm{PPL}}_N]=\sqrt{2\pi\sigma^2}\,(1-r/N)^{-N/2}
=\mathrm{PPL}\cdot e^{-r/2}(1-r/N)^{-N/2}.
\]
For $N\le r$ the expectation is infinite.

Expand $\log(1-r/N)=-r/N-r^2/(2N^2)+O(N^{-3})$:
\[
(1-r/N)^{-N/2}=e^{r/2}\big(1+\tfrac{r^2}{4N}+O(N^{-2})\big).
\]
So $c=r^2/4$. Independently, $\mathrm{Var}(x^2)=2s^4$ and $\mathrm{Var}(\ell)=r^2/2$, hence $\mathrm{Var}(\ell)/2=r^2/4$, matching problem~2.2(a).

If $r>1$, the model variance is smaller than the data variance. Then $\E[\widehat{\mathrm{PPL}}_N]=\infty$ for every $N\le r$: rare large data points make $-\log p_\theta(x)$ so heavy-tailed that the exponential moment does not exist. Consistency is only visible once $N>r$.
\end{hwsolbox}
